\chapter*{Conclusiones}
\addcontentsline{toc}{chapter}{Conclusiones}
\markboth{Conclusiones}{Conclusiones}

\lettrine[lines=2]{E}{n} este trabajo se desarrolló un formalismo vectorial para
calcular el campo electromagnético en el plano focal de un sistema estigmático
simple, formado por una única superficie cartesiana que separa dos medios
homogéneos. El formalismo combina la geometría de la superficie, descrita
por su radio de curvatura y la propiedad estigmática, la refracción
vectorial dada por los coeficientes de Fresnel y la propagación de
Fresnel hasta el plano focal, y produce el campo focal como transformada de
Fourier de una función pupila vectorial generalizada
$\mathbb E^{\mathscr G'}=\mathbb T\,\mathbb E_0/|\ell_0|$; bajo iluminación
azimutalmente simétrica esta transformada se reduce a transformadas de Hankel
de órdenes cero, uno y dos. Esta construcción se distingue de los formalismos
comúnmente empleados en difracción vectorial porque incorpora, punto a punto,
la transformación vectorial inducida por la superficie refractora, determinada
conjuntamente por su geometría y por su transmisión local de Fresnel.

El resultado central de este trabajo es que aun en los sistemas libres de
aberraciones geométricas están presentes, además de la aberración por
difracción, aberraciones de origen vectorial. La condición estigmática
elimina la aberración de camino óptico, asegurando fase constante en la
esfera convergente al punto imagen, pero al cambiar de medio las condiciones
de frontera de Maxwell~\cite{jackson1999} inducen sobre el campo una transformación vectorial
anisótropa y distinta en cada punto de la superficie, transformación que
localmente está caracterizada por los coeficientes de Fresnel. Sobre el plano transversal, esta transformación
es la de un diatenuador orientado radialmente en cada punto de la pupila,
cuyos modos propios son las polarizaciones radial y azimutal, moduladas
respectivamente por $\tau_p=t_p\cos\alpha'$ y $t_s$, donde $\cos\alpha'$
proyecta la componente $p$ refractada sobre el plano transversal. Las
combinaciones $t_\pm=\tau_p\pm t_s$
separan los dos efectos de la refracción: $t_+$ es una apodización común a
ambas componentes, mientras que la aberración vectorial, definida como la
desviación del campo focal respecto del límite escalar $t_s=\tau_p=1$ (en el
cual el formalismo recupera exactamente la teoría escalar de la difracción),
aparece siempre modulada por el factor $t_-$, no uniforme y no nulo fuera
del eje para una fuente puntual a distancia finita; en el límite de incidencia
colimada $t_-$ se anula en toda la pupila, de modo que los elipsoides, que
realizan ese límite, son superficies de interés para diseñar sistemas que
minimicen la aberración vectorial. En analogía con la teoría
de aberraciones de frente de onda, la aberración queda así caracterizada
sobre la propia pupila por la función matricial $\mathbb T$.

Las manifestaciones observables dependen del estado de polarización
incidente. En la base cilíndrica la transmisión es diagonal: las
polarizaciones radial y azimutal no sufren mezcla, y bajo iluminación
azimutalmente simétrica la estructura de Hankel de orden uno anula el campo
transversal en el eje, produciendo un anillo con centro oscuro. En las bases
que no coinciden con los modos propios, $t_-$ acopla polarizaciones
ortogonales: en la circular genera un canal
$\ket{L}\!\leftrightarrow\!\ket{R}$ con fase azimutal $e^{\pm2i\varphi}$,
suprimido sobre el eje donde $t_-$ y $J_2$ se anulan simultáneamente; en la
cartesiana genera una componente cruzada de orden angular dos, visible como
cuatro lóbulos en la polarización ortogonal a la incidente, que sobre el
estado de polarización actúa como una rotación local de la orientación sin
introducir elipticidad. Distintas polarizaciones incidentes producen así
distintos patrones de intensidad focal, algo que la teoría escalar, cuya
predicción es única e independiente de la polarización, no puede describir.

El régimen vectorial modifica también la forma del punto focal y, con ella,
la formación de imágenes. Para incidencia lineal la interferencia entre los
canales $\mathscr H_0$ y $\mathscr H_2$ rompe la simetría circular del lóbulo
central de la intensidad transversal: medido por su primer mínimo, el punto
queda caracterizado por dos radios principales que rotan rígidamente con la
polarización; para conjugados reales el punto se contrae a lo largo de ella y
se expande en la dirección perpendicular ($R_\parallel<R_\perp$). En el orden metaxial los
radios admiten una expresión cerrada; con conjugados reales, pupila uniforme
y cualquier vidrio óptico la apodización $t_+$ domina y ambos superan el radio de Airy,
$R_{\mathrm{esc}}<R_\parallel<R_\perp$. La deformación crece con la
desviación que impone la superficie; en el orden metaxial no depende del
índice, pero está acotada por $(n_i-n_0)/(n_i+n_0)$; la incidencia circular restaura la
simetría del punto por promedio azimutal. En la imagen de objetos extensos,
como mostró el ejemplo litográfico, la diatenuación de Fresnel hace que los
detalles alineados con la polarización incidente pierdan más contraste que
los ortogonales, y la iluminación circular reparte la pérdida entre ambas
orientaciones sin eliminarla. Se trata de una aberración de la misma
naturaleza que la que obliga a la litografía de alta apertura numérica a
controlar la polarización de la iluminación~\cite{flagello1996, ruoff2009},
presente aun cuando el sistema está libre de aberración geométrica.

El formalismo admite extensiones naturales: la reflexión se trata con $r_p,r_s$ y la geometría propia de la reflexión, en la que la cancelación del caso colimado no se produce; los sistemas de varias superficies se
tratan por composición de operadores, como ilustra el ejemplo litográfico. La
formulación del operador de transferencia como un operador espectral puede
usarse para generalizar el cálculo de difracción vectorial a iluminaciones
arbitrarias y a superficies dióptricas no necesariamente estigmáticas. La
verificación experimental de los resultados predichos, en particular la
componente cruzada de cuatro lóbulos y la deformación anisótropa del punto
focal bajo iluminación uniformemente lineal, serían extensiones naturales de este trabajo.
